\section{Limitations and Future Work}
\label{sec:limitations_future_work}

This work challenges a common assumption: that generative models require external, domain-specific encoders to improve representations and generation quality. As we show, external alignment can exhibit unexpected scaling behavior and often struggles to generalize across modalities. For example, REPA degrades audio generation compared to vanilla flow matching (Sec.~\ref{sec:experiments}). Instead, we address representation deficiency at its source by unifying generation and representation learning within a single framework.
This approach has trade-offs: the additional forward pass through the teacher increases training overhead, but the accelerated convergence and improved performance justify this cost (Fig.~\ref{fig:scaling}). Consistent with related work~\citep{sd3, rae}, the noise scheduler $p(t)$ requires tuning as it determines masking behavior. Fig.~\ref{fig:limitations} demonstrates this on text-to-image generation in the same setup described in Sec.~\ref{subsec:singlemodality}, where a uniform scheduler outperforms a logit-normal scheduler with a shift of $\alpha=1.78$ (App.~\ref{suppsubsec:aets}). While the better noise scheduling choice benefits both REPA and Self-Flow, the gap increases significantly in favor of the latter, which can be attributed to a more optimal timestep selection for the masking mechanism. In practice, we observe that the optimal scheduler for flow matching works well, and further tuning will likely yield additional gains.

Looking ahead, by bridging representation learning and generative modeling, our approach offers a path toward world models that harness the scalability and perceptual grounding of visual generative models without sacrificing the semantic abstraction required for planning and understanding - a direction we began to validate in Sec.~\ref{subsec:multimodal}. We hope this work stimulates research into consolidating generative and representation learning: two directions long pursued in isolation that may complement each other better than assumed.

